\section{Arquitetura do sistema}\label{sec:arquitetura}

A arquitetura do sistema foi definida com o objetivo de organizar as responsabilidades da aplicação e facilitar a manutenção, evolução e integração entre seus diferentes componentes. Considerando as características do sistema desenvolvido, adotou-se uma arquitetura composta por três camadas principais, apresentação, aplicação e persistência de dados, segundo o padrão de arquitetura em camadas discutido na Seção~\ref{sec:fund_arquitetura}. Essa organização permite separar a interface utilizada pelo profissional, as regras responsáveis pelo processamento das operações e o armazenamento das informações.

A camada de apresentação corresponde ao frontend da aplicação, desenvolvido utilizando React.js. É responsável pela interação com o usuário e pela disponibilização das funcionalidades do sistema por meio de uma interface web. Nessa camada são apresentadas as telas de autenticação, cadastro e consulta de aprendentes, gerenciamento dos instrumentos avaliativos, aplicação dos instrumentos e elaboração dos relatórios, destinadas ao psicopedagogo, e as de gerenciamento de usuários e consulta ao registro de auditoria, destinadas ao administrador. Em produção, o frontend é servido pelo Nginx, que também encaminha à API as requisições dirigidas a ela.

A comunicação entre a camada de apresentação e as funcionalidades do sistema ocorre por meio de uma API REST desenvolvida em Node.js utilizando o framework Express.js. A API funciona como intermediária entre o frontend e o banco de dados, recebendo as requisições realizadas pela interface, realizando as validações e regras de negócio correspondentes e retornando as informações necessárias para a aplicação. Seguindo o estilo REST \cite{fielding2000_rest}, a API não mantém estado de sessão no servidor: cada requisição carrega o token de acesso do usuário, e as operações sobre os recursos são expressas pelos métodos e códigos de estado do HTTP.

A camada de aplicação concentra as regras de negócio do sistema. Antes de chegar aos módulos de negócio, cada requisição passa pela autenticação, que identifica o usuário a partir do token de acesso, pela autorização por perfil e pela validação dos dados recebidos. Os módulos tratam da autenticação (\textit{auth}), dos usuários (\textit{users}), dos aprendentes (\textit{aprendentes}), dos instrumentos avaliativos (\textit{avaliacoes}), das aplicações desses instrumentos (\textit{aplicacoes}), dos relatórios (\textit{relatorios}), do painel geral (\textit{dashboard}) e do registro de auditoria (\textit{auditoria}). A utilização de uma camada específica para essas regras permite evitar que a lógica de negócio fique diretamente vinculada à interface, favorecendo a organização e a manutenção do código.

A persistência dos dados é realizada por meio de um banco de dados relacional MySQL. Essa camada é responsável pelo armazenamento estruturado das informações utilizadas pelo sistema, incluindo os dados dos usuários, dos aprendentes e de seu histórico, dos instrumentos e de suas perguntas, das aplicações e respostas, dos relatórios e do registro de auditoria. O modelo relacional permite estabelecer os vínculos entre essas entidades e preservar a integridade das informações.
Essa separação também contribui para a implementação dos mecanismos de segurança do sistema. A autenticação e a autorização são processadas no backend, permitindo que o acesso às operações da API seja controlado independentemente da interface apresentada ao usuário. Da mesma forma, o registro de auditoria é tratado no backend, na fronteira entre a camada de aplicação e a de persistência: as operações sobre o banco de dados passam por uma extensão do Prisma que registra quem a realizou, quando e sobre qual dado, sem depender da interface nem de cada regra de negócio.

A arquitetura adotada pode ser representada de forma simplificada pela Figura~\ref{fig:arquitetura_geral_do_sistema}, na qual são apresentadas as relações entre o frontend, a API e o banco de dados. A figura indica também a separação entre os perfis: os módulos em destaque reúnem dados clínicos e são exclusivos do psicopedagogo responsável, enquanto o administrador atua na gestão de usuários e na auditoria, sem acesso a esses dados (Seção~\ref{sec:separacao_admin}).

\begin{figure}[H]
    \centering
    \caption{Arquitetura geral do sistema}
    \label{fig:arquitetura_geral_do_sistema}
    \includegraphics[width=\textwidth,height=0.7\textheight,keepaspectratio]{arquitetura_sistema.png}
    \fonte{Elaborado pela autora (2026).}
\end{figure}
